\documentclass[12pt]{article}
\usepackage{fontspec}
\usepackage{xeCJK}
\setCJKmainfont[AutoFakeBold=2]{SimSun}
%\setCJKsansfont{SimHei}
\usepackage{amsmath, amssymb, amsthm}
\usepackage{geometry}
\usepackage{tikz}
\usepackage{fancyhdr}
%\setlength{\headheight}{15pt}
\usetikzlibrary{calc, intersections} % ABX 用於某些符號，若無可刪除
\usepackage{enumitem}
\usepackage{tkz-euclide}
\usepackage{subcaption}
\usepackage{indentfirst}
\usepackage{float}
\usepackage{chngcntr}
\counterwithin{figure}{section}
\counterwithin{figure}{subsection}
\usepackage[most]{tcolorbox}
\usepackage{etoolbox}
\usepackage{xcolor}
% 1. 定義一個新的開關，名字可以自訂（例如：showanswer）
\newtoggle{showanswer}

% 2. 設置開關狀態：true 為開啟（編譯），false 為關閉（不編譯）
\settoggle{showanswer}{true}

\newtcolorbox{mybox}[4][證明]{
  empty,
  breakable,                % 支援跨頁
  sharp corners,
  left=#2+#3,               % 為左側標籤預留空間 (縮進)
  right=0pt, top=0pt, bottom=0pt,
  colback=white,
  % --- 加入以下這行來調整段距 ---
  before upper={\setlength{\parskip}{0.5em plus 0.2em minus 0.1em}},
  after={\vspace{#4}},
  % ---------------------------
  overlay unbroken and first={ % 在第一頁/未斷頁時繪製標籤
    \node[
      anchor=north west,
      draw, 
      inner sep=0pt,
      minimum width=#2,
      minimum height=1.6em, 
      font=\bfseries\boldmath
    ] at (frame.north west) {#1};
  }
}

\geometry{left=1.5cm, right=1.5cm, top=2cm, bottom=2cm}
\counterwithout{figure}{subsection}
% 設定頁首頁尾
\pagestyle{fancy}
\fancyhf{}
\renewcommand{\headrulewidth}{2pt}
\renewcommand{\thefigure}{\thesection-\arabic{figure}} % 主圖編號
\renewcommand{\thesubfigure}{圖 \thesection-\arabic{figure}-\arabic{subfigure}} % 子圖變為 2-1
\captionsetup[figure]{labelformat=simple, labelsep=quad} 
\captionsetup[subfigure]{labelformat=simple, labelsep=quad} % 移除括號
\renewcommand{\figurename}{圖}
\lhead{廖汶鋒}
\chead{IMO 2026 Day 1}
\rhead{清華大學}
\fancyfoot[C]{\thepage} % Centered footer
%\fancyfoot[L]{2026年2月7日} % Left footer

%\title{四點共圓}
%\author{廖汶鋒}
%\date{2026年2月7日}
\linespread{1.3}

% 定段落首行縮進（兩個中文字）
\setlength{\parindent}{2em}
\setlength{\headheight}{15pt} % 避免頁首高度警告
\begin{document}
\thispagestyle{fancy}  % 讓首頁也顯示頁首
%\maketitle
\begin{center}
    \textbf{\Large{IMO 2026 Day 1}}
\end{center}

\begin{mybox}[第一題]{4em}{1em}{0.5em}
    黑板上寫有2026個大於1的正整數，允許有相同的數。在一次操作中，孔夫子選取黑板上兩個不同位置的整數 $m>1$ 和 $n>1$，
    並將這兩個數替換成
    $$  
        \gcd\left(m,n\right)\ \text{和}\ \frac{\text{lcm}\left(m,n\right)}{\gcd\left(m,n\right)}
    $$
    只要能夠進行這樣的操作他便繼續操作。
    \begin{enumerate}[label=(\arabic*)]
        \item 證明：不論孔夫子如何選擇，在有限次操作後，黑板上恰好有一個大於1的整數，記此數為$M$；
        \item 證明：整數$M$的值不依賴於孔夫子的選擇。
    \end{enumerate}
    （注：$\gcd\left(x,y\right)$表示正整數$x$和$y$的最大公因數，$\text{lcm}\left(x,y\right)$表示$x$和$y$的最小公倍數。）
\end{mybox}
\begin{mybox}[解答]{4em}{1em}{0.5em}
    \begin{enumerate}[label=(\arabic*)]
        \item 如果存在一種選擇，使得這樣的操作能夠無窮地進行下去，則等價於黑板上一直存在至少2種大於1的正整數。
        
        設第$k$輪操作後，黑板上大於1的正整數數量為$n_k$，同時記$n_0=2026$。結合假設，易知$\left\{n_k\right\}_{0}^{\infty}$是一個不增的無窮正整數序列，下界為2。
        因此，存在一個正整數$N$，使得對於所有的正整數$k\geq N$，都有$n_k=n_N$，即第$N$輪操作後，黑板上大於1的正整數數量不再變化，且不少於2。

        設在第$\ell$輪操作內，被選擇操作的兩個大於1的正整數為$m_{\ell}$和$n_{\ell}$；第$\ell$輪操作後，黑板上所有正整數的乘積為$A_{\ell}$。

        在第$\ell$輪操作後，被選擇操作的兩個數被替換為$\gcd\left(m_{\ell},n_{\ell}\right)$和$\frac{\text{lcm}\left(m_{\ell},n_{\ell}\right)}{\gcd\left(m_{\ell},n_{\ell}\right)}$，因此有
        \begin{align*}
            &\frac{A_{\ell+1}}{\gcd\left(m_{\ell},n_{\ell}\right)\cdot\frac{\text{lcm}\left(m_{\ell},n_{\ell}\right)}{\gcd\left(m_{\ell},n_{\ell}\right)}}=\frac{A_{\ell}}{m_{\ell}\cdot n_{\ell}}\\
            \Leftrightarrow\ &A_{\ell+1}=A_{\ell}\cdot\frac{\text{lcm}\left(m_{\ell},n_{\ell}\right)}{m_{\ell}\cdot n_{\ell}}=\frac{A_{\ell}}{\gcd\left(m_{\ell},n_{\ell}\right)}
        \end{align*}
        當$\ell\geq N$時，黑板上大於1的正整數數量不再變化，因此$\gcd\left(m_{\ell},n_{\ell}\right)\geq 2$，從而有$A_{\ell+1}\leq \frac{A_{\ell}}{2}$。

        另一方面，因為$n_N\geq 2$，所以對於任意正整數$k\geq N$，都有$A_k\geq 2^{n_N}\geq 2^2=4$。

        但是，當$\ell>\log_{2}{A_N}+N-2$時，則有
        $$
            4\leq A_{\ell}\leq\frac{A_N}{2^{\ell-N}}<4
        $$
        產生矛盾。因此，不論孔夫子如何選擇，題中所述的操作只能在有限次後結束，即黑板上存在至多一個大於1的正整數。

        如果在第$K$輪操作後，黑板上所有正整數都等於1，則可知
        $$
            \gcd\left(m_K,n_K\right)=\text{lcm}\left(m_K,n_K\right)=1
        $$
        即$m_K=n_K=1$，與題中所述操作的前提（$m_K,n_K>1$）矛盾。因此，黑板上恰好有一個大於1的正整數，證畢。

        \item 
    \end{enumerate}
\end{mybox}
\end{document}